% Copyright (c) 2026 Kitware, Inc.
% Author/provenance: Jon Crall, GPT 5.6 High
% Independent formalization-oriented literature distillation; not a transcription.
\documentclass{article}
\usepackage{amsmath,amssymb,mathtools}
\newcommand{\Hh}{\mathcal H}
\newcommand{\spec}{\operatorname{spec}}
\newcommand{\dist}{\operatorname{dist}}
\newcommand{\Ran}{\operatorname{Ran}}
\newcommand{\Graph}{\operatorname{Graph}}
\title{Infinite-dimensional Davis--Kahan theory: formalization distillation}
\author{Jon Crall and GPT 5.6 High}
\begin{document}
\maketitle

\section{Purpose and source discipline}
This document is an independent, formalization-oriented distillation of the
operator-theoretic arguments needed beyond the repository's finite-dimensional
Davis--Kahan development.  It records theorem interfaces, hidden analytic
lemmas, and proof decompositions.  It intentionally does not reproduce source
text.

Primary sources guiding the scaffold are Davis--Kahan (1970),
Kostrykin--Makarov--Motovilov on Riccati equations and generalized
$\tan 2\Theta$, Albeverio--Motovilov on the a priori $\tan\Theta$ theorem,
Seelmann on generic $\sin 2\Theta$ estimates, and the standard Kato theory of
closed operators and quadratic forms.

\section{Bounded Hilbert-space setup}
Let $A$ and $B$ be bounded self-adjoint operators on a Hilbert space $\Hh$.
A closed subspace $U$ reduces $A$ when both $U$ and $U^\perp$ are invariant.
Write $P_U$ for its orthogonal projection.  The gap metric is
$\|P_U-P_V\|$, while the directed gap is $\|P_{V^\perp}P_U\|$.

\section{Spectral separation}
The finite eigenvalue lists must be replaced by the spectrum of restrictions
to reducing subspaces.  The basic hypothesis is
\[
  \dist(\spec(A|_U),\spec(B|_{V^\perp}))\ge d>0.
\]
Convex-hull or interval/exterior separation gives the sharp Sylvester constant
one; arbitrary separated real spectra give the universal $\pi/2$ constant.

\section{Projection-valued spectral calculus}
For a bounded self-adjoint operator, the roadmap needs a Borel projection
$E_A(\Delta)$ satisfying orthogonality, multiplicativity, complementation,
strong countable additivity, commutation with $A$, and support on
$\spec(A)$.  Indicator functions in the bounded Borel calculus recover these
projections.  This is the main infrastructure gap between elementary compact
spectral theory and the general Hilbert-space theorem.

\section{Strong operator additivity}
Countable disjoint unions are represented by strong, not operator-norm,
convergence:
\[
  \sum_{j<n}E_A(\Delta_j)x\longrightarrow E_A(\cup_j\Delta_j)x.
\]
Lean therefore needs strong-operator convergence lemmas for projection sums,
orthogonal ranges, and passage through bounded operators.

\section{Functional calculus proof obligations}
The Borel calculus must be linear, multiplicative, star preserving, norm
controlled on the spectrum, and compatible with continuous functional
calculus.  Monotone approximation of indicators is the bridge to spectral
projections.

\section{Two-projection calculus}
The operator angle $\Theta(U,V)$ is recovered from the positive contraction
built from $P_U$ and $P_V$.  Its sine has norm $\|P_U-P_V\|$.  The acute case
$\|P_U-P_V\|<1$ is equivalent to bounded graph representability.

\section{Graph representation}
If $V$ is acute with respect to $U$, there is a unique bounded angular
operator $X:U\to U^\perp$ with $V=\Graph(X)$.  The identities
$\|X\|=\tan\|\Theta\|$ and the explicit block formula for $P_V$ provide the
bridge from subspace geometry to Riccati equations.

\section{Symmetric norm ideals}
In infinite dimension, general unitarily invariant norms are defined on
compact operator ideals.  The formal API must support two-sided ideal bounds,
pinching, unitary invariance, direct sums, singular-value gauges, and concrete
Schatten, trace-class, Hilbert--Schmidt, and Ky Fan instances.

\section{Sylvester equation}
For self-adjoint $A$ and $B$, solve
\[
  AX-XB=C.
\]
Under spectral separation the solution is unique.  Ordered or gap separation
yields $d\|X\|\le\|C\|$; arbitrary separated spectra yield
$d\|X\|\le(\pi/2)\|C\|$.  The proof uses a resolvent contour or Fourier/semigroup
integral and requires Bochner integration in operator spaces.

\section{Resolvent integral infrastructure}
Formalization requires resolvent existence away from the spectrum, norm
bounds by inverse spectral distance, measurability and integrability of the
operator-valued integrand, differentiation under the integral, and evaluation
of the Sylvester defect.

\section{Residual $\sin\Theta$ theorem}
For an isometric embedding $X$, a coordinate operator $M$, and residual
$R=AX-XM$, the cross block $P_{U^\perp}X$ satisfies a Sylvester equation.
Spectral separation gives
\[
 d\|P_{U^\perp}X\|\le\|R\|.
\]
This is the natural infinite-dimensional numerical-analysis interface.

\section{Perturbation $\sin\Theta$ theorem}
Taking $X$ to parametrize a reducing spectral subspace of $B$ converts the
residual to $(A-B)X$.  One mixed gap controls a directed angle.  Two mixed gaps
control the symmetric projector difference.  General spectral separation
inherits the $\pi/2$ Sylvester constant.

\section{$\sin 2\Theta$ theorem}
Reflection in a reducing subspace transforms the subspace problem into a
Sylvester equation for the off-diagonal projection block.  The reflection
defect is bounded by twice the perturbation, producing
\[
 d\|\sin(2\Theta)\|\le2\|A-B\|.
\]

\section{Generic double-angle bounds}
Without interval/exterior geometry, one combines the general Sylvester bound
with the two-projection calculus.  The resulting arcsine bounds must track the
branch of the maximal angle and the continuity path of spectral projections.

\section{Graph projections}
The projection onto $\Graph(X)$ is expressed using
$(I+X^*X)^{-1}$ and $(I+XX^*)^{-1}$.  Formalization needs positivity,
invertibility, square roots, and block multiplication identities.

\section{Riccati equation}
For the block operator
\[
 L=\begin{pmatrix}A_0&B\\B^*&A_1\end{pmatrix},
\]
$\Graph(X)$ is reducing exactly when
\[
 A_1X-XA_0-XBX+B^*=0.
\]
This equivalence is the central algebraic seam.

\section{Existence and uniqueness}
If $\spec(A_0)$ lies in a gap of $\spec(A_1)$ and the coupling is sufficiently
small, the Riccati equation has a bounded solution.  Contractive solutions
associated with the target spectral graph are unique.  Fixed-point,
continuation, and spectral-projection constructions should all be related.

\section{Block diagonalization}
A bounded Riccati solution gives a graph transform that block diagonalizes the
operator.  For self-adjoint data, the normalized graph transform is unitary.
This yields spectral enclosures and identifies the perturbed components.

\section{Off-diagonal gap preservation}
For an off-diagonal bounded self-adjoint perturbation $V$, the original gap
persists when $\|V\|<\sqrt2d$.  The proof combines spectral enclosure with
Riccati solvability and continuation of the spectral projection.

\section{Generalized $\tan 2\Theta$}
The angular operator of the perturbed reducing graph satisfies a sharp
quadratic inequality.  Translating its norm gives a two-sided estimate for
$\|P-Q\|$ and the classical $\tan 2\Theta$ upper bound.

\section{A priori $\tan\Theta$ theorem}
When one spectral component lies in a finite gap of the other, the sharp bound
is
\[
 \|P-Q\|\le \sin\!\left(\arctan\frac{\|V\|}{d}\right),
 \qquad \|V\|<\sqrt2d.
\]
The proof uses the Riccati solution and a refined spectral enclosure; the
threshold and bound are both sharp.

\section{Spectral repulsion}
Off-diagonal perturbations push the continued spectral components away from
the original gap.  The formal statement should be expressed both as set
inclusions and as lower bounds on the distance between restricted spectra.

\section{Direct rotation}
For acute subspaces there is a canonical unitary direct rotation intertwining
the projections.  Its square is the product of reflections.  The extremal
properties require careful angle hypotheses and should not be generalized to
all ideal gauges without proof.

\section{Closed operators}
Unbounded extensions require densely defined closed operators, graph norms,
extension relations, adjoints, resolvents, and spectral projections.  The
local Lean scaffold uses an explicit `ClosedOperator` structure while marking
reconciliation with `LinearPMap` as a foundational task.

\section{Kato--Rellich perturbations}
A bounded self-adjoint perturbation preserves self-adjointness on the original
domain.  More generally, a symmetric relatively bounded perturbation with
relative bound below one does so.  These results justify applying spectral
projection perturbation theory to unbounded operators.

\section{Unbounded spectral subspaces}
For self-adjoint closed operators, Borel spectral projections remain bounded
orthogonal projections.  If the operator difference is bounded, the bounded
Davis--Kahan arguments can often be recovered through resolvent identities.

\section{Unbounded Riccati equations}
Diagonal entries may be unbounded while off-diagonal blocks are bounded or
relatively bounded.  Strong solutions must preserve domains.  Reducing graph
subspaces are equivalent to strong Riccati solutions under explicit domain
regularity assumptions.

\section{Quadratic forms}
Form sums cover perturbations that are not operator bounded.  The KLMN theorem
constructs the perturbed self-adjoint operator.  A form-level Davis--Kahan
bound must measure the perturbation in a relative form norm and track how the
spectral gap changes.

\section{Form-domain geometry}
Proofs require Hilbert structures on form domains, bounded embeddings into the
ambient space, representation theorems, and resolvent estimates derived from
coercive forms.

\section{Compact spectral subspaces}
For compact self-adjoint operators, isolated nonzero spectral subspaces are
finite dimensional, but the ambient space need not be.  This is the first
practical bridge from the current repository to the full theory.

\section{Schatten consequences}
If the perturbation lies in a Schatten ideal, then under a gap hypothesis the
sine-angle operator and often the projector difference lie in the same ideal,
with the corresponding ideal-norm estimate.

\section{Hermitian dilation and Wedin}
For $T:\Hh_0\to\Hh_1$, the Hermitian dilation
$\begin{psmallmatrix}0&T^*\\T&0\end{psmallmatrix}$ converts singular subspaces
to spectral subspaces.  Applying the self-adjoint theory yields
infinite-dimensional Wedin bounds for isolated singular spectral sets.

\section{Sharpness}
All universal constants are tested on two-dimensional invariant blocks,
embedded into an arbitrary Hilbert space.  Orthogonal direct sums yield
simultaneous ideal-norm extremizers when the resulting operator belongs to the
chosen ideal.

\section{References}
\begin{itemize}
\item C. Davis and W. M. Kahan, ``The rotation of eigenvectors by a perturbation. III,'' SIAM J. Numer. Anal. 7 (1970), DOI 10.1137/0707001.
\item V. Kostrykin, K. A. Makarov, and A. K. Motovilov, ``On the existence of solutions to the operator Riccati equation and the $\tan\Theta$ theorem,'' arXiv:math/0210032.
\item V. Kostrykin, K. A. Makarov, and A. K. Motovilov, ``A generalization of the $\tan 2\Theta$ theorem,'' arXiv:math/0302020.
\item S. Albeverio and A. K. Motovilov, ``The a priori $\tan\Theta$ theorem for spectral subspaces,'' arXiv:1012.1569.
\item A. Seelmann, ``Notes on the $\sin 2\Theta$ theorem,'' Integral Equations Operator Theory (2014), DOI 10.1007/s00020-014-2127-z, arXiv:1310.2036.
\item K. A. Makarov, S. Schmitz, and A. Seelmann, ``Reducing graph subspaces and strong solutions to operator Riccati equations,'' arXiv:1307.6439.
\item T. Kato, Perturbation Theory for Linear Operators.
\end{itemize}
\end{document}
